\section{Passing to the limit}
\label{sec:lim}

Let $f_n \to f$ on a set $E$. Each property of the preceding sections may now
be asked twice: of the $f_n$, and of $f$. Figure~\ref{fig:blockc}
tabulates the verdicts; the rest of the section proves them.

\begin{figure}[htbp]
\centering
\begin{tabular}{@{}lcc@{}}
\toprule
if every $f_n$ has\ldots & \multicolumn{2}{c}{\ldots does the limit $f$?} \\
\cmidrule(l){2-3}
 & pointwise conv. & uniform conv. \\
\midrule
continuity                          & no (E9)  & \textbf{yes} \\
integrability                       & no (E7)  & \textbf{yes} \\
$\int_a^b f = \lim_n \int_a^b f_n$  & no (E11) & \textbf{yes} \\
\addlinespace[3pt]
differentiability                   & no       & \textbf{no} (E8, E10) \\
\bottomrule
\end{tabular}
\caption{What survives a limit. Uniform convergence is strictly stronger
than pointwise convergence, and it rescues the first three rows; the fourth
it cannot rescue, because differentiation multiplies amplitude by frequency
and uniformity bounds only amplitude. The repair for the last row is
Theorem~\ref{thm:repair}: assume uniform convergence of the derivatives
$f_n'$ instead.}
\label{fig:blockc}
\end{figure}

Recall Definition~\ref{def:conv}: convergence is \emph{pointwise} if
$f_n(x)\to f(x)$ at each $x$, and \emph{uniform} if
$\sup_E|f_n-f|\to0$, so that one index serves every point at once.
Figure~\ref{fig:tube} shows the gap between the two on $\mathrm{E}9$.

\begin{figure}[htbp]
\centering
\begin{tikzpicture}[x=1mm, y=1mm]
  \fill[black!9] (0,0) rectangle (52,5);
  \draw[axis] (0,0) -- (60,0);
  \draw[axis] (0,0) -- (0,34);
  \draw[guide] (0,26) -- (51,26);
  \node[mlbl, anchor=east] at (-1.4,26) {$1$};
  \draw[thin, dash pattern=on 2pt off 2pt] (0,5) -- (56,5);
  \node[mlbl, anchor=east] at (-1.4,5) {$\varepsilon$};
  \foreach \p in {1.4,4,9}{
    \draw[thin] plot[smooth, domain=0:52, samples=40]
      (\x, {26*(\x/52)^\p});}
  \node[mlbl, anchor=south east] at (50,20) {$x^{n}$};
  \draw[seg] (0,0) -- (51,0);
  \node[ptopen] at (52,0) {};
  \node[ptfill] at (52,26) {};
  \node[mlbl, anchor=west] at (53.6,26) {$f(1)=1$};
  \node[mlbl, anchor=north] at (52,-1.6) {$1$};
  \node[lbl, anchor=west] at (2,31)
    {each $x^n$ climbs out of the tube before reaching $1$};
  \draw[thin, -{Stealth[length=3pt]}] (14,29) -- (43.5,5.8);
  \node[lbl, anchor=west, align=left] at (58,8)
    {the $\varepsilon$-tube about\\ the limit function};
  \draw[thin, -{Stealth[length=3pt]}] (57.4,6.6) -- (49,2.5);
  \node[lbl, anchor=west] at (10,-7)
    {the limit $f$ (heavy): $0$ on $[0,1)$, $1$ at $x=1$};
  \draw[thin, -{Stealth[length=3pt]}] (22,-5.6) -- (30,-0.7);
\end{tikzpicture}
\caption{$f_n(x)=x^n$ on $[0,1]$ ($\mathrm{E}9$; three members drawn, later
ones lower). The pointwise limit $f$ --- the heavy graph, $0$ up to $1$ and
then the isolated value $1$ --- is discontinuous. Uniform convergence would
require the whole graph of $x^n$, for large $n$, to stay inside the shaded
$\varepsilon$-tube about $f$; however large $n$ is, the graph still climbs
out before reaching $1$, so the convergence is pointwise but not uniform.}
\label{fig:tube}
\end{figure}

\begin{theorem}\label{thm:unifcont}
If each $f_n$ is continuous and $f_n \to f$ uniformly, then $f$ is
continuous.
\end{theorem}

\begin{flow}
  \item \textit{Locate the obstacle.} We must compare $f(x)$ with $f(x_0)$,
        and we know nothing about $f$ directly --- it is only a limit.
  \item \textit{Find a detour.} Each $f_n$ \emph{is} continuous, so route
        through one: from $f(x)$ up to $f_n(x)$, across to $f_n(x_0)$, and
        back down to $f(x_0)$ (Figure~\ref{fig:3eps}).
  \item \textit{Price the legs.} The two vertical legs are short by
        convergence; the horizontal leg is short by continuity of $f_n$.
        Three legs, so each is allotted $\varepsilon/3$.
  \item \textit{Check the order of choices.} This is where uniformity is
        spent. The index $n$ must be fixed \emph{before} $x$ is, because the
        $\delta$ of the horizontal leg depends on $n$. Uniform convergence
        permits that order; pointwise convergence does not, since there the
        admissible $n$ depends on the point.
  \item \textit{Audit.} Uniformity is used exactly twice, for the two
        vertical legs, and continuity of a single $f_n$ once. No property of
        the index beyond its existence is used, which is why the same
        argument proves the corresponding statements for uniform limits of
        uniformly continuous or of bounded functions.
\end{flow}

\begin{proof}
Fix $x_0$ and $\varepsilon>0$. Choose $n$ with
$\sup|f_n-f|<\varepsilon/3$, then $\delta$ with
$|f_n(x)-f_n(x_0)|<\varepsilon/3$ for $|x-x_0|<\delta$. For such $x$,
\[ |f(x)-f(x_0)| \le |f(x)-f_n(x)|+|f_n(x)-f_n(x_0)|+|f_n(x_0)-f(x_0)|
   < \varepsilon. \qedhere \]
\end{proof}

\begin{figure}[htbp]
\centering
\begin{tikzpicture}[x=1mm, y=1mm]
  \node[ptfill] (a) at (0,0)   {};
  \node[ptfill] (b) at (0,22)  {};
  \node[ptfill] (c) at (58,22) {};
  \node[ptfill] (d) at (58,0)  {};
  \node[mlbl, anchor=east]  at (-2,0)  {$f(x)$};
  \node[mlbl, anchor=east]  at (-2,22) {$f_n(x)$};
  \node[mlbl, anchor=west]  at (60,22) {$f_n(x_0)$};
  \node[mlbl, anchor=west]  at (60,0)  {$f(x_0)$};
  \draw[rmarrow] (a) -- node[lbl, anchor=east, inner sep=3pt]
    {$<\varepsilon/3$ (uniform)} (b);
  \draw[rmarrow] (b) -- node[lbl, anchor=south, inner sep=3pt]
    {$<\varepsilon/3$ ($f_n$ continuous)} (c);
  \draw[rmarrow] (c) -- node[lbl, anchor=west, inner sep=3pt]
    {$<\varepsilon/3$ (uniform)} (d);
  \draw[guide] (a) -- node[lbl, anchor=north, inner sep=3pt]
    {what we must bound: $<\varepsilon$} (d);
\end{tikzpicture}
\caption{The three legs of the $3\varepsilon$ argument. The direct route
along the bottom is unavailable, since nothing is known about $f$; one
detours through a single $f_n$. Uniform convergence is what permits that $n$
to be chosen \emph{before} $x$, which is exactly what pointwise convergence
would not allow.}
\label{fig:3eps}
\end{figure}

Uniformity is used for the outer two terms. Under pointwise convergence each
can still be made small, but only after the point has been fixed, and the
middle term then requires a $\delta$ depending on the index chosen for that
point; the estimate does not close. $\mathrm{E}9$ is what the failure looks
like.

\begin{theorem}\label{thm:unifint}
If each $f_n$ is Riemann integrable on $[a,b]$ and $f_n\to f$ uniformly, then
$f$ is integrable and $\int_a^b f = \lim_n \int_a^b f_n$.
\end{theorem}

\idea{Uniform convergence puts $f$ inside a horizontal band of width
$2\varepsilon_n$ about $f_n$. Upper and lower sums are monotone under such a
comparison, so the strips of Figure~\ref{fig:sums} for $f$ are no taller than
those for $f_n$ plus the band --- and the band contributes at most
$2\varepsilon_n(b-a)$ no matter how the partition is chosen. Two smallnesses
are combined: one from $n$, one from the partition.}

\begin{proof}
Put $\varepsilon_n = \sup_{[a,b]}|f_n-f| \to 0$, so that
$f_n-\varepsilon_n \le f \le f_n+\varepsilon_n$. For any partition $P$,
\[ U(P,f)-L(P,f) \le U(P,f_n)-L(P,f_n)+2\varepsilon_n(b-a), \]
and choosing $n$ and then $P$ suitably makes the right side small, giving
integrability. The same inequalities bound
$\bigl|\int f - \int f_n\bigr|$ by $\varepsilon_n(b-a)$.
\end{proof}

\begin{theorem}\label{thm:unifdiff}
Uniform convergence does not preserve differentiability; and even when the
limit is differentiable, $\lim_n f_n'$ need not equal $f'$.
\end{theorem}

\begin{proof}
For the first clause, take the partial sums of $\mathrm{E}8$: each is a
finite sum of cosines, hence $\mathscr{C}^\infty$, and by the $M$-test they
converge uniformly --- to Weierstrass's function, which is differentiable at
no point. So a uniform limit of smooth functions need not be differentiable
anywhere.

For the second, let $f_n(x)=\sin(nx)/\sqrt n$ on $\R$
($\mathrm{E}10$; Figure~\ref{fig:wiggle}).
Then $\sup|f_n| = 1/\sqrt n \to 0$, so $f_n \to 0$ uniformly and the limit
is smooth. But $f_n'(x)=\sqrt n\cos(nx)$ is unbounded and has no limit at
any point; at $x=0$, $f_n'(0)=\sqrt n \to \infty$ while $f'\equiv 0$. So
even when the limit is differentiable, the derivatives may converge to
nothing at all.
\end{proof}

\begin{figure}[htbp]
\centering
\begin{tikzpicture}[x=1mm, y=1mm]
  \draw[axis] (0,14) -- (46,14);
  \draw[thin, dash pattern=on 2pt off 2pt] (0,17.7) -- (44,17.7);
  \draw[thin, dash pattern=on 2pt off 2pt] (0,10.3) -- (44,10.3);
  \node[mlbl, anchor=west] at (45,17.7) {$1/\sqrt n$};
  \foreach \p/\amp in {2/9, 5/5.7}{
    \draw[thin] plot[smooth, domain=0:44, samples=120]
      (\x, {14 + \amp*sin(\p*\x*4)});}
  \draw[seg] plot[smooth, domain=0:44, samples=120]
    (\x, {14 + 3.7*sin(12*\x*4)});
  \node[lbl, anchor=north] at (22,-4) {$f_n = \sin(nx)/\sqrt n \to 0$};
  \begin{scope}[xshift=62mm]
  \draw[axis] (0,14) -- (46,14);
  \draw[thin, dash pattern=on 2pt off 2pt] (0,26) -- (44,26);
  \draw[thin, dash pattern=on 2pt off 2pt] (0,2) -- (44,2);
  \node[mlbl, anchor=west] at (45,26) {$\sqrt n$};
  \foreach \p/\amp in {2/5, 5/8}{
    \draw[thin] plot[smooth, domain=0:44, samples=120]
      (\x, {14 + \amp*cos(\p*\x*4)});}
  \draw[seg] plot[smooth, domain=0:44, samples=120]
    (\x, {14 + 12*cos(12*\x*4)});
  \node[lbl, anchor=north] at (22,-4) {$f_n' = \sqrt n\,\cos(nx)$ diverges};
  \end{scope}
\end{tikzpicture}
\caption{Why uniform convergence cannot control derivatives. The same three
functions appear in both panels, and the heavy curve is the largest $n$ of
the three. On the left it is the \emph{flattest}: its amplitude $1/\sqrt n$
(dashed) has collapsed furthest. On the right, after differentiating, the
very same curve is the \emph{steepest}: the factor $n$ overwhelms the
collapse, and its amplitude $\sqrt n$ (dashed) has grown furthest. Uniform
smallness of a function says nothing about its slope: a wiggle may be short
and steep at once.}
\label{fig:wiggle}
\end{figure}

\begin{theorem}[Rudin 7.17]\label{thm:repair}
Suppose each $f_n$ is differentiable on $[a,b]$, the sequence
$\bigl(f_n(x_0)\bigr)$ converges for some $x_0\in[a,b]$, and $(f_n')$
converges uniformly. Then $(f_n)$ converges uniformly to some $f$, and
$f'=\lim_n f_n'$.
\end{theorem}

The hypothesis has moved onto the derivatives, and the functions themselves
need only converge at one point. Differentiation is not continuous with
respect to uniform convergence, so to differentiate a limit one must assume
something about the derivatives; the same asymmetry is why term-by-term
differentiation of a power series requires its own theorem while term-by-term
integration does not.

\begin{theorem}\label{thm:ptwise}
Pointwise convergence preserves neither continuity, nor integrability, nor
the value of the integral.
\end{theorem}

\begin{proof}
Continuity fails by $\mathrm{E}9$. For integrability, enumerate
$\Q\cap[0,1]=\{q_1,q_2,\dots\}$ and let $f_n = \1_{\{q_1,\dots,q_n\}}$; each
$f_n$ has finitely many discontinuities, hence is integrable by
Theorem~\ref{thm:lebesgue}, while $f_n \to \mathrm{E}7$ pointwise and
$\mathrm{E}7$ is not integrable. For the integral, the spikes
$\mathrm{E}11$ (Figure~\ref{fig:spikes}) satisfy $f_n \to 0$ pointwise with $\int_0^1 f_n = 1$ for
every $n$.
\end{proof}

\begin{figure}[htbp]
\centering
\begin{tikzpicture}[x=1mm, y=1mm]
  \fill[black!7]  (0,0) rectangle (24,8);
  \fill[black!13] (0,0) rectangle (12,16);
  \fill[black!19] (0,0) rectangle (6,32);
  \draw[axis] (0,0) -- (56,0);
  \draw[axis] (0,0) -- (0,36);
  \foreach \w/\h/\lab in {24/8/{n=1}, 12/16/{n=2}, 6/32/{n=4}}{
    \draw[thin] (0,\h) -- (\w,\h) -- (\w,0);
    \node[mlbl, anchor=west] at (\w+1,\h) {$\lab$};}
  \draw[thin, dash pattern=on 2pt off 2pt] plot[smooth]
    coordinates {(24,8) (12,16) (8,24) (6,32)};
  \node[mlbl, anchor=east] at (-1.4,8)  {$1$};
  \node[mlbl, anchor=east] at (-1.4,16) {$2$};
  \node[mlbl, anchor=east] at (-1.4,32) {$4$};
  \node[mlbl, anchor=north] at (6,-1.6)  {$\tfrac14$};
  \node[mlbl, anchor=north] at (12,-1.6) {$\tfrac12$};
  \node[mlbl, anchor=north] at (24,-1.6) {$1$};
  \node[lbl, anchor=west, align=left] at (26,26)
    {the corners ride the hyperbola $xy=1$:\\
     every spike has area $n \times \tfrac1n = 1$};
  \draw[thin, -{Stealth[length=3pt]}] (25.4,24.2) -- (9.5,23);
  \node[lbl, anchor=west] at (8,-7)
    {at each fixed $x>0$ the spikes eventually lie to its left, so
     $f_n(x)\to 0$};
\end{tikzpicture}
\caption{The spikes $f_n = n\1_{(0,1/n)}$ ($\mathrm{E}11$), each enclosing
area $1$: as the base shrinks to $1/n$ the height grows to $n$, and the
top-right corner $(1/n,\,n)$ rides up the hyperbola $xy=1$. At each fixed
$x>0$ the spike eventually lies entirely to the left of $x$, so
$f_n(x)\to0$ pointwise; yet $\int f_n = 1$ for every $n$. The mass escapes
upward, which no pointwise statement detects.}
\label{fig:spikes}
\end{figure}

\begin{remark}
Pointwise limits are nonetheless not arbitrary. A pointwise limit of
continuous functions is of Baire class one, and not every function is: the
Dirichlet function $\mathrm{E}7$ is not a pointwise limit of continuous
functions, by Baire's theorem \cite[Lect.~28]{yupin}. It does arise as a
pointwise limit of the integrable functions in the proof above, but never of
continuous ones.
\end{remark}

\begin{remark}
The uniformity hypothesis in Theorem~\ref{thm:unifint} is a limitation of the
Riemann integral rather than of the statement. Under the Lebesgue integral
the conclusion $\int \lim = \lim \int$ holds under a pointwise hypothesis
together with an integrable dominating function --- a hypothesis the spikes
of $\mathrm{E}11$ conspicuously fail.
\end{remark}
